\documentclass[conference]{IEEEtran}
\IEEEoverridecommandlockouts

\usepackage{cite}
\usepackage{amsmath,amssymb,amsfonts}
\usepackage{algorithmic}
\usepackage{algorithm}
\usepackage{graphicx}
\usepackage{textcomp}
\usepackage{xcolor}
\usepackage{booktabs}
\usepackage{multirow}
\usepackage{subcaption}

\begin{document}

\title{EXPA: Multi-Stage Adversarial Patch Defense Through \\
Cascade Detection and Explainable AI Localization}

\author{\IEEEauthorblockN{Hussein Harissa}
\IEEEauthorblockA{\textit{Computer Science and Networks Department} \\
\textit{ENSTA Bretagne}\\
Brest, France \\
hussein.harissa@ensta-bretagne.fr}
}

\maketitle
\begin{abstract}
Adversarial patch attacks, demonstrated by Brown et al. to fool deep
neural networks through physically realizable perturbations, pose
serious security threats to image classification systems. We validate
this threat by generating universal adversarial patches via Projected
Gradient Descent against AlexNet (ImageNet pretrained), achieving
attack success rates ranging from 4.9\% to 51.7\% depending on target
class --- revealing that patch effectiveness is strongly governed by the
topology of the target class in the model's feature space.

To defend against such attacks, we present EXPA, a four-stage
preprocessing-based defense pipeline. Our cascade architecture operates
as follows: (1) a median filter KL-divergence screen flags suspicious
images by measuring prediction shift under spatial smoothing; (2)
EigenCAM activation mapping confirms the attack by detecting abnormally
concentrated class-relevant regions in the feature space; (3) EigenCAM
heatmap thresholding precisely localizes the patch bounding box; and
(4) median filter inpainting reconstructs the patch region to recover
the original class prediction. This dual use of the median filter ---
for both detection and recovery --- makes the pipeline fully
gradient-free, requiring no model retraining or backpropagation.

Evaluated on ImageNette across three target classes (banana, bee, lemon),
EXPA achieves up to 99.8\% attack detection with 19.6\% false positive
rate and 100\% patch localization accuracy. Class restoration via median
inpainting averages 58.4\% across all patch/dataset combinations, with
66.1\% on the banana patch. Results demonstrate that gradient-free
explainability tools like EigenCAM can be effectively repurposed from
model interpretation to robust adversarial defense through strategic
cascade design.
\end{abstract}


\begin{IEEEkeywords}
Adversarial machine learning, explainable AI, adversarial patches, EigenCAM, image classification security, cascade defense
\end{IEEEkeywords}

\section{Introduction}

\subsection{Context}

Deep neural networks (DNNs) have achieved remarkable success in computer vision, surpassing human-level performance on numerous benchmarks. However, they exhibit a critical vulnerability: susceptibility to adversarial examples---inputs crafted to cause misclassification \cite{szegedy2013intriguing}. Among adversarial attacks, \textbf{adversarial patches} \cite{brown2017adversarial} represent a particularly dangerous class because they are physically realizable: attackers can print patches on physical media and place them in real-world environments to cause persistent misclassification.

\begin{figure}[t]
\centering
\includegraphics[width=0.48\textwidth]{figures/patch_attack_examples.pdf}
\caption{Adversarial patch attacks across domains: (a) stop sign fooled by sticker, (b) facial recognition defeated by eyeglass patch, (c) person invisible to surveillance \cite{eykholt2018robust,sharif2016accessorize,thys2019fooling}.}
\label{fig:patch_examples}
\end{figure}

The threat has been demonstrated across critical domains: autonomous vehicles misclassifying stop signs \cite{eykholt2018robust}, facial recognition systems fooled by wearable patches \cite{sharif2016accessorize}, surveillance cameras bypassed through adversarial clothing \cite{thys2019fooling}, and medical imaging systems producing incorrect diagnoses \cite{finlayson2019adversarial}. What makes patches particularly dangerous is their \textbf{universality}: a single optimized design causes misclassification across diverse images, viewing angles, and lighting conditions, enabling mass production and deployment at scale.

We validate this threat by generating universal patches via Projected Gradient Descent (PGD) \cite{madry2018towards} against AlexNet \cite{krizhevsky2012imagenet}. Across multiple target classes (banana, bee, lemon) and datasets (ImageNette validation and training sets), we observe attack success rates ranging from 4.9\% to 51.7\%. This variation reveals that patch effectiveness is strongly governed by target class topology in feature space---isolated classes are more vulnerable while classes with many feature-space neighbors prove more resistant.

\subsection{Problem Statement}

Defending against adversarial patches requires simultaneously addressing four challenges:

\textbf{(1) High Detection Accuracy:} Reliably identify patches (target: $>$90\% detection rate).

\textbf{(2) Low False Positive Rate:} Avoid excessive false alarms (target: $<$10\% false positives).

\textbf{(3) Computational Efficiency:} Execute within real-time constraints for practical deployment.

\textbf{(4) No Model Modification:} Work with existing trained models without retraining.

Existing defenses face fundamental trade-offs:

\begin{itemize}
\item \textbf{Adversarial training} \cite{madry2018towards} improves robustness but requires extensive retraining (days to weeks), degrades clean accuracy by 5--10\%, and fails to generalize to novel attack types.

\item \textbf{Certified defenses} \cite{cohen2019certified} provide provable guarantees but incur prohibitive computational costs ($>$1 second per image), making real-time deployment infeasible.

\item \textbf{Input preprocessing} \cite{guo2017countering,xu2017feature} is computationally efficient but suffers high false positives ($>$30\%) as natural textures also shift predictions after filtering.

\item \textbf{XAI-based detection} like SentiNet \cite{chou2020sentinet} achieves $\sim$92\% detection but relies on heavy segmentation models adding $\sim$200ms overhead, exceeding real-time constraints.
\end{itemize}

\begin{figure}[t]
\centering
\includegraphics[width=0.48\textwidth]{figures/defense_tradeoffs.pdf}
\caption{Existing defenses face trade-offs between accuracy, speed, and false positives. Our cascade approach achieves practical balance: 96.4\% detection with 19.6\% false positive rate on the ImageNette banana patch.}
\label{fig:tradeoffs}
\end{figure}

No existing approach balances all four requirements for practical deployment.

\subsection{Summary of the Solution}

We present a four-stage cascade defense leveraging EigenCAM \cite{muhammad2021eigen}---a gradient-free explainable AI technique---to detect, localize, and mitigate adversarial patches. The key insight: patches must dominate model attention to cause misclassification, creating abnormally concentrated activation patterns detectable through explainability analysis.

\begin{figure*}[t]
\centering
\includegraphics[width=0.95\textwidth]{figures/pipeline_overview.pdf}
\caption{Four-stage defense pipeline: (a) Median filter detection via KL divergence (patches: KL$>$1.0, clean: KL$<$0.5), (b) EigenCAM confirmation using four anomaly metrics, (c) Heatmap thresholding for localization (100\% accuracy), (d) Median inpainting for recovery (58.4\% average class restoration across all patches).}
\label{fig:pipeline}
\end{figure*}

Our pipeline consists of four sequential stages:

\textbf{Stage 1 --- Median Filter Detection.} Apply $5\times5$ median filter and measure KL divergence between original and filtered predictions. Patches are fragile to smoothing (KL$>$1.0), clean images stable (KL$<$0.5). Provides fast initial screening with high sensitivity ($>$99\% TPR) but elevated false positives ($\sim$34\%).

\textbf{Stage 2 --- EigenCAM Confirmation.} For flagged images, extract gradient-free EigenCAM heatmaps from early (features.3: $192\times27\times27$) and late (features.10: $256\times13\times13$) AlexNet layers. Compute four anomaly metrics---spatial entropy, peak/mean ratio, top-$k$ energy, cross-layer correlation---and use majority voting to confirm attacks. Reduces false positives by 43\% relative to Stage 1 alone (34.5\% $\rightarrow$ 19.6\%).

\textbf{Stage 3 --- Patch Localization.} Threshold EigenCAM heatmap at top-5\% highest-activation pixels. Connected component analysis yields tight bounding box. Achieves 100\% localization accuracy on ImageNette because patches inherently create concentrated attention patterns.

\textbf{Stage 4 --- Inpainting and Recovery.} Replace localized patch region with median-filtered content. Re-classify recovered image. Success measured by target class removal ($\sim$100\%) and original class restoration (58.4\% average across all patches). Dual use of median filtering (Stages 1 and 4) makes pipeline fully gradient-free.

\section{Related Work}

\subsection{Adversarial Patch Attacks}

Brown et al. \cite{brown2017adversarial} introduced adversarial patches as physically realizable adversarial examples. Unlike $\ell_p$-bounded perturbations affecting all pixels minimally, patch attacks apply arbitrary perturbations to a localized region. The optimization objective maximizes target class probability:

\begin{equation}
\delta^* = \arg\max_{\delta} \mathbb{E}_{x \sim \mathcal{D}} \left[ P(y_t | x \oplus \delta) \right]
\end{equation}

where $\delta$ is the patch, $x$ is the input image, $\oplus$ denotes patch application, and $y_t$ is the target class.

Physical implementations have demonstrated real-world viability: Eykholt et al. \cite{eykholt2018robust} attacked traffic sign classifiers using printed stickers, achieving high success rates on stop sign misclassification. Sharif et al. \cite{sharif2016accessorize} created wearable patches that fool facial recognition systems, enabling impersonation attacks against state-of-the-art face recognition. Thys et al. \cite{thys2019fooling} developed adversarial patches that, when worn on clothing, make people invisible to person detection systems. These works collectively demonstrate that patches are not theoretical curiosities but practical security threats that can be deployed in real-world environments with minimal cost.

\subsection{Explainable AI Techniques}

Class Activation Mapping (CAM) \cite{zhou2016learning} pioneered neural network visualization by projecting final layer weights onto convolutional feature maps, highlighting discriminative regions for specific classes. Grad-CAM \cite{selvaraju2017grad} generalized this approach to arbitrary CNN architectures using gradient-based weighting:

\begin{equation}
\alpha_k = \frac{1}{Z} \sum_i \sum_j \frac{\partial y^c}{\partial A^k_{ij}}
\end{equation}

where $\alpha_k$ are importance weights for channel $k$, $y^c$ is the class score, $A^k$ are activation maps, and $Z$ is the spatial dimension. The final heatmap combines weighted activations:

\begin{equation}
L_{Grad-CAM} = ReLU\left(\sum_k \alpha_k A^k\right)
\end{equation}

While widely adopted for model interpretation, Grad-CAM requires backpropagation through the network to compute gradients, adding computational overhead and potential numerical instability.

\textbf{EigenCAM} \cite{muhammad2021eigen} eliminates gradient computation through Principal Component Analysis (PCA) on activation maps. Muhammad and Yeasin report approximately 40\% speedup compared to Grad-CAM while maintaining comparable visualization quality. Given activation maps $A \in \mathbb{R}^{C \times H \times W}$ reshaped to $A' \in \mathbb{R}^{(HW) \times C}$ (spatial positions $\times$ channels), EigenCAM centers the data, computes the covariance matrix, and extracts the first principal component via eigendecomposition:

\begin{equation}
L_{EigenCAM} = \text{reshape}(A' \cdot v_1)
\end{equation}

where $v_1$ is the first eigenvector. The projection onto this component yields a spatial heatmap indicating which image regions most strongly influence the model's prediction. The gradient-free nature makes EigenCAM particularly suitable for applications where backpropagation overhead is prohibitive.

\subsection{XAI-Based Defenses}

SentiNet \cite{chou2020sentinet} pioneered the application of explainability tools to adversarial patch detection. The key observation is that adversarial patches create abnormally concentrated Grad-CAM activation patterns: to cause misclassification, patches must dominate the model's attention, leading to spatially localized high-activation regions distinct from normal image features. SentiNet's three-stage pipeline consists of: (1) Grad-CAM-based detection measuring attention concentration, (2) Mask R-CNN instance segmentation for precise patch localization, and (3) inpainting-based removal and re-classification. The authors report approximately 92\% detection rate on their evaluation dataset, though the system incurs substantial computational overhead (approximately 200ms per image) due to the Mask R-CNN segmentation model, which limits applicability to real-time scenarios.

Beyond SentiNet, other XAI-based approaches have been explored. Ma et al. \cite{ma2018characterizing} use Local Intrinsic Dimensionality (LID) to detect adversarial examples through feature space analysis, though their method was primarily designed for $\ell_p$-bounded perturbations rather than localized patches. These works collectively demonstrate that explainability tools can reveal adversarial perturbations through analysis of model activations and feature representations.

\subsection{Preprocessing-Based Defenses}

Input transformation defenses modify images before classification to neutralize adversarial perturbations. Xu et al. \cite{xu2017feature} introduced Feature Squeezing, which combines bit-depth reduction and spatial smoothing. Guo et al. \cite{guo2017countering} systematically evaluated multiple transformation techniques including JPEG compression, total variation minimization, and image quilting.

\textbf{Median filtering} is particularly effective against adversarial patches due to its non-linear, outlier-rejecting properties. For a $5\times5$ window, the median is the 13th value when sorted---robust to up to 12 outliers per window. Unlike Gaussian blur, median filtering isolates and eliminates adversarial pixel values without averaging them into surrounding pixels.

However, preprocessing defenses used in isolation produce high false positive rates. Natural images with complex textures (brick walls, fabric patterns) shift model predictions after filtering, triggering false alarms. This fundamental trade-off motivates our cascade design.

\subsection{Adversarial Training and Certified Defenses}

\textbf{Adversarial training} \cite{madry2018towards} augments training data with adversarial examples generated via PGD, forcing models to learn robust features. While effective, it requires extensive GPU computation, degrades clean accuracy by 5--10\%, and fails to generalize to novel attack types.

\textbf{Certified defenses} \cite{cohen2019certified} provide provable robustness guarantees through randomized smoothing or interval bound propagation. However, these methods incur prohibitive computational costs and have limited applicability to localized patch attacks, which violate the uniform perturbation assumptions underlying most certification techniques.

\section{Methodology}

\subsection{Experimental Setup}

\subsubsection{Model Architecture}

We evaluate our cascade defense pipeline on AlexNet \cite{krizhevsky2012imagenet}, an 8-layer convolutional neural network pretrained on ImageNet (1000 classes, 1.2M training images). Input resolution: $224\times224$ RGB images.

For EigenCAM analysis, we extract activation maps from two layers:
\begin{itemize}
\item \textbf{Early layer (features.3):} 192 channels, $27\times27$ spatial resolution. Captures low-level features such as edges, textures, and simple patterns.
\item \textbf{Late layer (features.10):} 256 channels, $13\times13$ spatial resolution. Captures high-level semantic features. Adversarial patches typically hijack these late layers.
\end{itemize}

\subsubsection{Datasets}

\textbf{ImageNette:} A 10-class subset of ImageNet (tench, English springer, cassette player, chain saw, church, French horn, garbage truck, gas pump, golf ball, parachute). The validation set contains 3,925 test images; the training set contains 9,469 images.

\subsubsection{Adversarial Patch Generation}

We generate universal adversarial patches using PGD \cite{madry2018towards}:

\begin{equation}
\delta_{t+1} = \Pi_{[0,1]} \left( \delta_t + \alpha \cdot \text{sign}\left(\nabla_\delta \mathcal{L}(f(x \oplus \delta_t), y_t)\right) \right)
\end{equation}

\textbf{Optimization Parameters:}
\begin{itemize}
\item Iterations: 100; Step size: $\alpha = 0.05$
\item Patch size: $28\times28$ pixels (1.56\% of $224\times224$ image area)
\item Patch position: Centered; Initialization: Random uniform $\in [0,1]$
\end{itemize}

\textbf{Target Classes:} Class 954 (banana), Class 309 (bee), Class 951 (lemon)---all semantically unrelated to ImageNette categories.

\subsection{Defense Pipeline Implementation}

\subsubsection{Stage 1 --- Median Filter Detection}

Apply a $5\times5$ median filter to obtain $\tilde{x}$, then compute KL divergence:

\begin{equation}
D_{KL}(f(x) \| f(\tilde{x})) = \sum_{i=1}^{1000} f(x)_i \log\frac{f(x)_i}{f(\tilde{x})_i}
\end{equation}

If $D_{KL} > \tau = 1.0$, flag the image as suspicious and pass to Stage 2.

\subsubsection{Stage 2 --- EigenCAM Confirmation}

Extract EigenCAM heatmaps from both layers. Compute four anomaly metrics on the late-layer heatmap:

\textbf{1. Spatial Entropy:} $H = -\sum_{i,j} p_{ij} \log p_{ij}$ (low entropy $\Rightarrow$ concentrated attention)

\textbf{2. Peak-to-Mean Ratio:} $R = \max(L)/\text{mean}(L)$ (high ratio $\Rightarrow$ localized spike)

\textbf{3. Top-$k$ Energy:} $E_k = \sum_{\text{top-5\%}} L_{ij} / \sum L_{ij}$ (high $\Rightarrow$ small region dominates)

\textbf{4. Cross-Layer Correlation:} $\rho = \text{corr}(L_{early}, L_{late})$ (low $\Rightarrow$ layer disagreement)

Majority voting ($\geq 2$ metrics flagged) plus a mask-confidence-drop test confirms attacks.

\subsubsection{Stage 3 --- Patch Localization}

Threshold heatmap at top-5\% highest-activation pixels, compute minimal bounding box with 8-pixel padding:

\begin{equation}
M = \mathbb{I}(L > \text{percentile}(L, 95))
\end{equation}

\subsubsection{Stage 4 --- Inpainting and Recovery}

Replace localized region with median-filtered content:

\begin{equation}
x'_{ij} = \begin{cases}
\tilde{x}_{ij} & \text{if } (i,j) \in [y_1', y_2'] \times [x_1', x_2'] \\
x_{ij} & \text{otherwise}
\end{cases}
\end{equation}

Re-classify $x'$ and measure: (1) target class removed, (2) original class restored.

\section{Experimental Results}

\subsection{Attack Success Rate Analysis}

Table~\ref{tab:attack_success} reports attack success rates across all target classes and datasets.

\begin{table}[h]
\centering
\caption{Attack success rates across target classes and datasets (28$\times$28 patch)}
\label{tab:attack_success}
\begin{tabular}{lcc}
\toprule
\textbf{Target Class} & \textbf{ImageNette val} & \textbf{ImageNette train} \\
\midrule
Class 954 (banana) & 193/3925 (4.9\%)  & 265/9469 (2.8\%) \\
Class 309 (bee)    & 618/3925 (15.7\%) & 1114/9469 (11.8\%) \\
Class 951 (lemon)  & 2029/3925 (51.7\%) & 844/9469 (8.9\%) \\
\midrule
\textbf{Range}     & \textbf{4.9\%--51.7\%} & \textbf{2.8\%--11.8\%} \\
\bottomrule
\end{tabular}
\end{table}

Attack success rates vary dramatically (4.9\% to 51.7\%) depending on target class selection, revealing that patch effectiveness is strongly governed by the target class's topology in AlexNet's feature space. The lemon patch achieves a strikingly high 51.7\% ASR on the validation set, likely because yellow-fruit classes (banana, lemon) occupy adjacent regions in AlexNet's feature space, creating a broad attractor basin that the patch easily exploits. Classes that are feature-space neighbors to many ImageNet training classes prove more resistant to patch attacks.

\subsection{Cascade Defense Performance}

Table~\ref{tab:cascade} compares cascade defense performance across datasets.

\begin{table}[h]
\centering
\caption{Cascade defense performance (28$\times$28 patch, class 954 ``banana'')}
\label{tab:cascade}
\begin{tabular}{lccc}
\toprule
\textbf{Dataset} & \textbf{Fooled} & \textbf{TPR (\%)} & \textbf{FPR (\%)} \\
\midrule
ImageNette val   & 193  & 96.4 & 19.6 \\
ImageNette train & 265  & 98.9 & 19.2 \\
\bottomrule
\end{tabular}
\end{table}

The cascade architecture achieves high detection rates ($>$96\%) while maintaining consistent false positive rates around 19\% on both val and train sets, confirming the reliability of the threshold.

\subsection{Cascade Components Comparison}

Table~\ref{tab:components} shows the contribution of each stage.

\begin{table}[h]
\centering
\caption{Defense component comparison (ImageNette val, 28$\times$28 patch, class 954, 3{,}925 clean images)}
\label{tab:components}
\begin{tabular}{lcc}
\toprule
\textbf{Defense Component} & \textbf{TPR (\%)} & \textbf{FPR (\%)} \\
\midrule
Median Filter (Stage 1 only)               & 99.5 & 34.5 \\
\textbf{Cascade (Median + EigenCAM)}       & \textbf{96.4} & \textbf{19.6} \\
\midrule
\textbf{FPR Reduction}                     & ---  & \textbf{43.2\%} \\
\bottomrule
\end{tabular}
\end{table}

The cascade achieves 43.2\% false positive reduction compared to median filtering alone (34.5\% $\rightarrow$ 19.6\%) while maintaining high detection sensitivity. This validates that EigenCAM confirmation effectively filters false alarms from natural image textures.

\subsection{End-to-End Pipeline Results --- All 6 Runs}

Table~\ref{tab:allruns} reports the complete four-stage pipeline performance across all patch/dataset combinations.

\begin{table}[h]
\centering
\caption{Module 10 end-to-end results: all target classes and datasets (28$\times$28 patch, top-5\% EigenCAM)}
\label{tab:allruns}
\resizebox{\columnwidth}{!}{%
\begin{tabular}{llrccccc}
\toprule
\textbf{Dataset} & \textbf{Patch} & \textbf{N} & \textbf{TPR} & \textbf{FPR} & \textbf{Loc.} & \textbf{Cls.~Rest.} & \textbf{E2E} \\
\midrule
val   & banana & 193  & 96.4\% & 19.6\% & 100\% & 66.1\% & 95.9\% \\
val   & bee    & 618  & 99.4\% & ---    & 100\% & 45.8\% & 99.4\% \\
val   & lemon  & 2029 & 99.8\% & ---    & 100\% & 54.0\% & 99.8\% \\
train & banana & 265  & 98.9\% & 19.2\% & 100\% & 66.4\% & 98.9\% \\
train & bee    & 1114 & 99.8\% & ---    & 100\% & 59.1\% & 99.8\% \\
train & lemon  & 844  & 97.9\% & ---    & 100\% & 59.2\% & 97.9\% \\
\midrule
\multicolumn{2}{l}{\textbf{Average}} & & \textbf{98.7\%} & \textbf{19.4\%} & \textbf{100\%} & \textbf{58.4\%} & \textbf{98.6\%} \\
\bottomrule
\end{tabular}%
}
\end{table}

\textit{Note: FPR (---) for bee/lemon patches shares the same clean set as banana; the measured FPR of 19.6\% (val) and 19.2\% (train) applies uniformly per dataset.}

\subsection{Impact of Localization Threshold}

\begin{table}[h]
\centering
\caption{EigenCAM top-$k$ threshold impact (ImageNette val, 28$\times$28, class 954)}
\label{tab:topk}
\begin{tabular}{lcc}
\toprule
\textbf{top-$k$} & \textbf{Localization (\%)} & \textbf{Class Restored (\%)} \\
\midrule
10\% & 100.0 & 54.8 \\
\textbf{5\%}  & \textbf{100.0} & \textbf{66.1} \\
\bottomrule
\end{tabular}
\end{table}

Tightening the threshold from 10\% to 5\% improves class restoration by 11.3 percentage points (54.8\% $\rightarrow$ 66.1\%) while maintaining perfect localization accuracy.

\section{Discussion}

\subsection{Detection Performance Analysis}

The cascade architecture successfully balances sensitivity and specificity. Median filter provides high sensitivity (Stage 1 TPR: 99.5\%--100\% across all runs) by flagging any significant prediction shift under spatial smoothing, while EigenCAM confirmation reduces false alarms by 43.2\% relative to Stage 1 alone (34.5\% $\rightarrow$ 19.6\%).

The missed detections (0.2\%--3.6\% depending on patch) correspond to patches where PGD optimization converged suboptimally, producing adversarial signals insufficiently fragile to median smoothing. Threshold tuning (lowering from 1.0 to 0.7--0.8) could improve detection sensitivity at the cost of increased false positives.

False positives (19.6\% on ImageNette val) arise from complex natural textures---brick walls, patterned fabrics, repetitive architectural elements---that exhibit both prediction shift under smoothing and moderate spatial attention concentration. The cascade's 43.2\% FPR reduction compared to median filtering alone demonstrates that EigenCAM confirmation successfully filters many texture-induced false alarms.

\subsection{Localization Accuracy}

Perfect 100\% localization accuracy across all six runs demonstrates a fundamental theoretical property: adversarial patches are inherently visible to explainability tools due to their attention-dominating mechanism. For a patch to cause misclassification to target class $y_t$, PGD optimization forces it to strongly activate feature detectors associated with $y_t$. This necessarily concentrates model attention on the patch region, making it the dominant direction in EigenCAM's PCA. The patch reveals itself through its very mechanism of attack.

\subsection{Class Restoration and Recovery Limitations}

Class restoration varies by patch type: banana 66.1\%/66.4\%, lemon 54.0\%/59.2\%, bee 45.8\%/59.1\% (val/train). Target class removal is near-perfect across all runs ($\geq$99.5\%), confirming that inpainting effectively neutralizes the adversarial signal. Restoration failures fall into two categories:

\textbf{1. Semantically Critical Patch Placement ($\sim$60\% of failures):} Patches covering discriminative image regions (object centers, distinctive features) cannot be reconstructed by simple median filtering, causing AlexNet to misclassify the recovered image.

\textbf{2. Ambiguous Post-Inpainting Features ($\sim$40\% of failures):} Median filtering produces smoothed content lacking sufficient high-frequency detail for confident classification, causing recovery to a plausible but incorrect class.

The bee patch achieves lowest class restoration on val (45.8\%) partly because bee images in ImageNet share visual features with other insects and animals, creating a more ambiguous recovery context. Future work could explore deep learning-based inpainting (LaMa \cite{suvorov2022lama}, MAT) for semantically-aware reconstruction.

\subsection{Feature Space Topology and Attack Vulnerability}

The dramatic variation in attack success rates (4.9\%--51.7\%) reveals a novel empirical finding: patch effectiveness is governed by the target class's position in AlexNet's feature space. The lemon patch achieves 51.7\% ASR because yellow-fruit classes (banana, lemon) are feature-space neighbors in AlexNet---a patch optimized for ``banana'' simultaneously activates ``lemon'' detectors, creating a broad attractor basin. In contrast, the banana patch achieves only 4.9\% ASR against the same ImageNette images, suggesting that the banana class requires more precise activation patterns that are harder to achieve universally.

\subsection{Limitations and Threats to Validity}

\textbf{1. False Positive Rate:} The 19.6\% FPR may be unacceptable for high-stakes deployments. Mitigation strategies include stricter thresholds, unanimous metric agreement, and temporal validation in video contexts.

\textbf{2. Non-Adaptive Attacks:} Our evaluation assumes black-box attackers unaware of the defense. Adaptive attackers could craft patches minimizing KL divergence (robust to median filtering) or distributing attention to evade EigenCAM detection. Evaluating adaptive attack robustness is critical future work.

\textbf{3. Single Architecture:} Results are evaluated on AlexNet. Generalization to modern architectures (ResNet, ViT) requires additional evaluation, though the EigenCAM-based approach is architecture-agnostic.

\textbf{4. Inpainting Quality:} Median filter inpainting is computationally efficient but produces lower visual quality than deep inpainting methods. The 58.4\% average class restoration rate demonstrates feasibility, with room for improvement through learned reconstruction.

\section{Conclusion}

We presented EXPA, a four-stage gradient-free defense against adversarial patch attacks. By combining median filter KL-divergence detection, EigenCAM-based confirmation, heatmap localization, and inpainting recovery, our cascade achieves:

\begin{itemize}
\item \textbf{98.7\% average TPR} across 6 patch/dataset combinations
\item \textbf{19.6\% cascade FPR} (43.2\% reduction vs.~Stage 1 alone)
\item \textbf{100\% patch localization} --- universal across all tested patches
\item \textbf{58.4\% class restoration} via simple median inpainting
\item \textbf{Fully gradient-free} --- no backpropagation or model modification required
\end{itemize}

Results across three target classes (banana, bee, lemon) and two dataset splits confirm system reliability. The 100\% localization accuracy across all configurations demonstrates that adversarial patches are inherently detectable via explainability analysis due to their attention-dominating mechanism.

The feature-space topology finding---that attack success rates vary dramatically (4.9\%--51.7\%) based on target class selection---opens new research directions in understanding adversarial vulnerability at the feature space level. Future work includes: adaptive attack evaluation, deep learning-based inpainting for higher class restoration, extension to modern architectures (ResNet, ViT), and real-time implementation optimization.

\bibliographystyle{IEEEtran}
\bibliography{references}

\end{document}
